\section{Introduction}

\paragraph{Background.}
Sign language is one of the most important communication methods~\cite{who2021}
used by deaf and hard-of-hearing individuals around the world. It lets people with
hearing disabilities communicate using hand gestures, facial expressions, and body
movements instead of spoken language. However, learning sign language can be hard
for beginners because there are few accessible educational tools and interactive
learning systems.

With the progress of artificial intelligence and computer vision, deep learning
models can now recognize complex visual patterns such as hand
gestures~\cite{krizhevsky2012}. This makes it possible to build systems that
recognize sign language automatically from camera input. Yet most existing systems
focus on American Sign Language (ASL)~\cite{rastgoo2021survey}, while Arabic Sign
Language (ArSL) gets far less attention~\cite{mohandes2014} and has fewer
interactive educational applications, even though Arabic is among the most spoken
languages in the world and its deaf community is large.

Beyond language coverage, research has mostly emphasized accuracy on benchmark
datasets, without always asking whether that accuracy reflects real performance. A
high number on a benchmark does not automatically make a useful tool: the score can
be inflated by how the data is split and evaluated, and lab conditions differ from
real-world use. Closing the gap between reported benchmarks and practical
educational tools therefore needs attention not only to interaction and real-time
performance, but also to whether the evaluation itself is valid.

\subsection{Problem Statement}
Although recent deep learning systems report strong performance in sign language
recognition, most work focuses on ASL, while ArSL still suffers from limited
datasets, fewer educational applications, and a lack of interactive real-time
systems built for Arabic-speaking users.

Many existing systems also work only as static classifiers that predict a gesture
without offering an engaging learning experience, with little support for
interactive practice or real-time feedback.

A further, often ignored problem is the validity of the reported accuracy itself.
Sign-language datasets are frequently built from many near-duplicate frames of the
same signer. When such data is split randomly at the image level, almost-identical
frames land in both the training and test sets. This leaks information and inflates
test accuracy, so the number no longer reflects how the model performs on new
users. A high score obtained this way can hide poor real-world performance.

Finally, real deployment brings variation in lighting, hand position, camera
quality, and background clutter that can lower performance outside the lab, which
is especially clear in webcam-based systems.

There is therefore a need for an Arabic Sign Language learning platform that
combines \emph{honestly evaluated} recognition with an interactive tutoring
experience that works under practical conditions.

\subsection{Project Objectives}
The main objective is to build an interactive real-time Arabic Sign Language
learning system using deep learning and computer vision. The specific objectives
are:
\begin{itemize}
  \item Develop a real-time Arabic sign language recognition system using a
        fine-tuned MobileNetV2 model.
  \item Recognize the Arabic alphabet sign classes using the ArASL dataset.
  \item Establish an honest, leakage-free evaluation, and measure generalization on
        an independent dataset instead of trusting a single benchmark score.
  \item Build an interactive tutoring interface that guides users through word
        spelling letter-by-letter with immediate feedback.
  \item Improve robustness to real-world conditions such as background clutter and
        lighting changes.
  \item Deploy the system with a lightweight architecture suitable for real-time
        webcam use.
\end{itemize}

\subsection{Project Contributions}
The main contributions are:
\begin{itemize}
  \item An interactive Arabic Sign Language learning system that combines deep
        learning recognition with educational tutoring.
  \item A real-time webcam-based recognition pipeline using a fine-tuned
        MobileNetV2 architecture.
  \item A word-based spelling tutor that guides users through sequential letter
        prediction and practice.
  \item Hand-focused preprocessing and background mitigation to improve robustness.
  \item \textbf{Identifying and quantifying data leakage} in the standard random
        split of frame-based ArSL datasets, and showing that a small
        train/validation gap under such a split is \emph{not} evidence of
        generalization.
  \item \textbf{A leakage-aware (cluster-based) split and cross-dataset evaluation
        protocol} that gives a realistic generalization estimate instead of an
        inflated benchmark score.
\end{itemize}